% !TEX root = ../main.tex
\chapter{鞅}
\label{ch:martingale}

\noindent\emph{用到的前置工具：条件期望及其塔性质（``条件期望''一章，作为投影理解）；概率空间、$\sigma$-代数与可测性（概率与测度部分）；$\sigma$-代数作为``信息''的解读与后验信念的鞅性（``贝叶斯更新与信息''一章）。}
\medskip

经济学里反复出现一句口头禅：``有效市场上没有免费的午餐''、``理性预期下的误差无法预测''、``公平的赌局长期下来不赚不赔''。这三句话听起来像格言，其实背后是同一个精确的数学对象——\textbf{鞅}（martingale）。鞅是``公平博弈''的形式化：给定到此刻为止的全部信息，下一刻的最优预测就是当前值，既不系统性地上漂、也不系统性地下沉。一旦把``不可被系统性预测''这件直觉的事钉成定义，有效市场假说、理性预期、贝叶斯学习、无套利定价就都成了同一句话的不同方言。

本章按使用者的顺序展开。先把承载``信息随时间积累''的 filtration 讲清楚（它把上一部分的 $\sigma$-代数串成一条时间线）；再用条件期望写出鞅的定义，并分出上鞅、下鞅与鞅差分序列这几个亲戚；接着给两条撑起全部应用的大定理——\emph{可选停时定理}（赌徒无法靠聪明的停手策略稳赚）与\emph{鞅收敛定理}（鞅在温和条件下必收敛），证明以讲清直觉为主、极重的部分只陈述；最后把工具对准它的四个去处：有效市场与风险中性定价、理性预期的预测误差、时间序列的鞅中心极限定理、以及序贯检验。读完应当能看出：鞅不是随机过程里一个孤立的技术概念，而是``信息''与``期望''这两件旧工具在时间轴上的自然合体。

\section{信息流：filtration}
\label{sec:martingale-filtration}

要谈``随时间到来的信息'',先得让``信息''有一个随时间变化的载体。上一部分我们已经学会用一个 $\sigma$-代数 $\cF$ 表示``此刻能分辨的事件全体''。现在把时间放进来：在每个时刻 $t$，我们掌握的信息对应一个 $\sigma$-代数 $\cF_t$，而``信息只增不减''——昨天知道的今天仍然知道——就要求这列 $\sigma$-代数\emph{递增}。

\defn{Filtration（信息流）}{
设 $(\Omega,\cF,\PP)$ 为概率空间，时间指标 $t=0,1,2,\dots$（或连续 $t\ge 0$）。一列子 $\sigma$-代数
\[
    \cF_0\se\cF_1\se\cF_2\se\cdots\se\cF
\]
若满足 $s\le t\Rightarrow\cF_s\se\cF_t$，则称为一个 \textbf{filtration}（信息流）。其中 $\cF_t$ 解读为``截止到时刻 $t$ 所能知道的一切信息''。
}

最常见的 filtration 是由一个随机过程\emph{自身}的历史生成的：给定过程 $\set{X_t}$，令
\[
    \cF_t=\sigma(X_0,X_1,\dots,X_t),
\]
即``观察到截止 $t$ 的全部取值''所携带的信息，称为该过程的\emph{自然 filtration}。一个过程若在每个时刻 $t$ 都是 $\cF_t$-可测的（即``到 $t$ 为止的信息足以算出 $X_t$''），就说它关于这个 filtration 是\textbf{适应的}（adapted）。直觉上，适应性排除了``偷看未来''：$X_t$ 的取值只能依赖到 $t$ 为止已经发生的事，不能依赖明天的随机性。

\state{filtration 把``信息一章''搬上了时间轴}{
上一部分讲``信息越多 $\Rightarrow$ $\sigma$-代数越细''，那是\emph{静态}的两个信息结构在比精粗。filtration 不过是把这件事\emph{动态化}：让信息结构沿时间 $t$ 一格一格地变细，$\cF_0\se\cF_1\se\cdots$ 就是``学习''在数学上的样子——每过一期，划分被进一步切碎、能分辨的事件变多。鞅、停时、布朗运动全都建在这条信息流上。每当你在宏观或金融模型里看到 $\E{\cdot\given\cF_t}$（``给定 $t$ 期信息取期望''），背后就是这一条 filtration。
}

有了 filtration，``给定到 $t$ 为止的信息，对未来取期望''就有了精确写法：$\E{Y\given\cF_t}$。这正是上一章条件期望的运算，只不过条件化的对象是一个随时间增长的信息流。鞅的定义，就是对这个量提一个出奇简单的要求。

\section{鞅、上鞅、下鞅}
\label{sec:martingale-def}

\defn{鞅、上鞅、下鞅}{
设 $\set{\cF_t}$ 为 filtration，过程 $\set{X_t}$ 关于它适应且每个 $X_t$ 可积（$\E{\abs{X_t}}<\infty$）。称 $\set{X_t}$ 为
\begin{itemize}
    \item \textbf{鞅}（martingale），若对一切 $t$，$\ \E{X_{t+1}\given\cF_t}=X_t$；
    \item \textbf{下鞅}（submartingale），若 $\ \E{X_{t+1}\given\cF_t}\ge X_t$；
    \item \textbf{上鞅}（supermartingale），若 $\ \E{X_{t+1}\given\cF_t}\le X_t$。
\end{itemize}
}

定义只有一行，却要逐字咀嚼。条件 $\E{X_{t+1}\given\cF_t}=X_t$ 说的是：\emph{站在今天、用今天的全部信息去预测明天，最优预测恰好就是今天的值}。鞅既不预期上涨、也不预期下跌——它是``公平''的精确含义。下鞅有系统性的上漂（``对你有利的博弈''，财富在期望意义下增长），上鞅则有系统性的下沉（``对你不利的博弈'')。图~\ref{fig:martingale-1} 把鞅画了出来：已实现的过去是一条确定的折线，而站在当前时刻 $t$ 往后看，所有可能的未来路径在条件期望的意义上恰好围绕着水平线 $X_t$ 展开。

\begin{figure}[h]
\centering
\begin{tikzpicture}[scale=1.0,>=Stealth,line cap=round,line join=round]
  % 鞅样本路径：已实现的过去 + 当前时刻 t 处未来的条件期望持平于 X_t。
  \draw[->] (-0.2,0) -- (10.6,0) node[right] {$s$};
  \draw[->] (0,-0.2) -- (0,4.6) node[above] {$X_s$};

  \def\tstar{5}
  \def\xstar{2.4}

  % 已实现过去路径（折线，至 t）
  \draw[very thick,blue!60!black]
      (0,2.0) -- (1,2.6) -- (2,2.1) -- (3,2.9) -- (4,2.2) -- (\tstar,\xstar);
  \fill[blue!60!black] (\tstar,\xstar) circle (2.4pt);

  % 持平水平线：未来条件期望 = X_t
  \draw[red!75!black,thick,dash pattern=on 5pt off 3pt]
      (\tstar,\xstar) -- (10.2,\xstar);

  % 四条可能的未来路径：两条上漂、两条下沉，终点关于持平线 X_t 上下均衡，
  % 视觉上恰好围绕 X_t 展开（终点均值落在红线上）。
  \draw[gray!65] (\tstar,\xstar) -- (6,2.7) -- (7,3.1) -- (8,2.6) -- (9,3.3) -- (10,3.5);
  \draw[gray!65] (\tstar,\xstar) -- (6,2.9) -- (7,2.4) -- (8,3.0) -- (9,2.6) -- (10,2.9);
  \draw[gray!65] (\tstar,\xstar) -- (6,1.9) -- (7,2.4) -- (8,1.8) -- (9,2.2) -- (10,1.9);
  \draw[gray!65] (\tstar,\xstar) -- (6,2.1) -- (7,1.6) -- (8,2.2) -- (9,1.5) -- (10,1.3);

  % 竖直虚线标记当前时刻
  \draw[dashed,gray] (\tstar,0) -- (\tstar,\xstar);
  \node[anchor=north] at (\tstar,-0.05) {$t$};

  % 当前值标签（轴外）
  \draw[dashed,gray] (\tstar,\xstar) -- (0,\xstar);
  \node[anchor=east] at (-0.1,\xstar) {$X_t$};

  % 持平线标签：停在右上空白，细引线指向红线
  \node[red!75!black,anchor=west,align=left] at (7.1,4.35)
       {$\E{X_{s}\given\cF_t}=X_t\ (s>t)$};
  \draw[red!75!black,->] (8.4,4.10) -- (8.85,{\xstar+0.07});

  \node[blue!60!black,anchor=south] at (2.3,3.15) {已实现};
  \node[gray!50!black,anchor=north] at (7.9,0.95) {可能的未来路径};
\end{tikzpicture}
\caption{鞅的样本路径。蓝色折线是已实现的过去（在 $\cF_t$ 下确定）；站在当前时刻 $t$ 向后看，未来充满随机性（几条灰色路径），但它们的条件期望恰好持平于当前值 $X_t$（红色水平虚线）。``下一步往哪走''无法预测，``下一步的平均位置''就是脚下——这就是公平。}
\label{fig:martingale-1}
\end{figure}

\rmkb[``给定当前信息''与``给定当前值''不是一回事]{
鞅定义里条件化的是整个信息流 $\cF_t$（直到 $t$ 的全部历史），不是单个当前值 $X_t$。若 $\E{X_{t+1}\given X_t}=X_t$（只用当前值）成立但 $\E{X_{t+1}\given\cF_t}=X_t$（用全部历史）不成立，那它就\emph{不}是鞅。这个区别在经济学里要紧：``价格是鞅''意味着\emph{任何}基于公开历史信息的交易策略都无法系统性获利，而不只是``不能靠当前价格预测''。Markov 性（只看当前状态）与鞅性（条件期望持平）是两条独立的性质——一个过程可以是其一而非其二。
}

把定义里的 $X_t$ 移到等号一侧，鞅的``增量''有一个极有用的名字。

\defn{鞅差分序列}{
设 $\set{X_t}$ 是关于 $\set{\cF_t}$ 的鞅，记增量 $D_t:=X_t-X_{t-1}$。则 $\set{D_t}$ 称为 \textbf{鞅差分序列}（martingale difference sequence, MDS），它满足
\[
    \E{D_{t+1}\given\cF_t}=\E{X_{t+1}-X_t\given\cF_t}=X_t-X_t=0 .
\]
即\emph{给定过去，每一步增量的条件期望为零}。反过来，$X_t=X_0+\sum_{s=1}^{t}D_s$ 把鞅写成``初值加一串鞅差分的累加''。
}

鞅差分是``不可预测的增量''的精确刻画。它比``独立同分布''弱、却往往够用：增量不必独立、不必同分布，只要``给定过去、当前增量的最优预测是零''。这一条性质有一个立竿见影的推论——\textbf{鞅差分两两不相关}：对 $s<t$，
\[
    \E{D_s D_t}=\E{\E{D_s D_t\given\cF_{t-1}}}=\E{D_s\,\E{D_t\given\cF_{t-1}}}=\E{D_s\cdot 0}=0 ,
\]
其中第二步用了``提出已知量''（$D_s$ 在 $\cF_{t-1}$ 下已知，可提到条件期望外，因为 $s<t$）、第三步用了鞅差分性质。这条``不相关''正是后面鞅中心极限定理和理性预期检验的支点。

\state{鞅差分序列：计量经济学里``扰动''的正确身份}{
本科计量常假设误差项 i.i.d.，但时间序列里这太强：今天的冲击常与昨天的波动相关。真正需要的、也几乎处处够用的，是把扰动设成\emph{鞅差分序列}——$\E{\varepsilon_{t+1}\given\cF_t}=0$。它恰好等价于``\emph{用过去信息无法预测下一期扰动}''，正是理性预期、有效市场、以及 GMM 矩条件 $\E{g(\cdot)\given\cF_t}=0$ 想表达的东西。鞅差分两两不相关、却允许条件异方差（如 ARCH/GARCH 的波动聚集），这使它成为金融时间序列的默认扰动结构。
}

下鞅与上鞅则刻画``有方向的漂移''，图~\ref{fig:martingale-2} 把三者并排对照：从同一个当前值出发，下一期的条件期望或上移（下鞅）、或持平（鞅）、或下移（上鞅）。

\begin{figure}[h]
\centering
\begin{tikzpicture}[scale=1.0,>=Stealth,line cap=round,line join=round]
  % 下鞅 / 鞅 / 上鞅：当前值 X_t 与下一期条件期望 E[X_{t+1}|F_t] 的关系。
  \draw[->] (-0.2,0) -- (11.0,0) node[right] {$t$};
  \draw[->] (0,-0.2) -- (0,4.7) node[above] {$X_t$};

  \def\xt{2.3}
  \def\tnow{3}
  \def\tnext{5}

  \draw[very thick,black!70]
      (0,1.9) -- (1,2.4) -- (2,2.0) -- (\tnow,\xt);
  \fill[black!70] (\tnow,\xt) circle (2.2pt);

  % 水平参考线只延到下一期结点即止，避免灰虚线穿过右侧蓝色“鞅”标签
  \draw[dashed,gray!60] (\tnow,\xt) -- (\tnext,\xt);
  \draw[dashed,gray!60] (\tnow,\xt) -- (0,\xt) node[anchor=east] {$X_t$};
  \draw[dashed,gray!55] (\tnow,0) -- (\tnow,\xt);
  \node[anchor=north] at (\tnow,-0.05) {$t$};
  \draw[dashed,gray!55] (\tnext,0) -- (\tnext,3.55);
  \node[anchor=north] at (\tnext,-0.05) {$t+1$};

  \def\subY{3.35}
  \draw[->,very thick,green!50!black] (\tnow,\xt) -- (\tnext,\subY);
  \fill[green!50!black] (\tnext,\subY) circle (2.0pt);

  \def\marY{2.3}
  \draw[->,very thick,blue!60!black] (\tnow,\xt) -- (\tnext,\marY);
  \fill[blue!60!black] (\tnext,\marY) circle (2.0pt);

  \def\supY{1.25}
  \draw[->,very thick,red!70!black] (\tnow,\xt) -- (\tnext,\supY);
  \fill[red!70!black] (\tnext,\supY) circle (2.0pt);

  \node[green!50!black,anchor=west] at (\tnext+0.15,\subY)
       {下鞅：$\E{X_{t+1}\given\cF_t}\ge X_t$};
  \node[blue!60!black,anchor=west] at (\tnext+0.15,\marY)
       {鞅：$\E{X_{t+1}\given\cF_t}=X_t$};
  \node[red!70!black,anchor=west] at (\tnext+0.15,\supY)
       {上鞅：$\E{X_{t+1}\given\cF_t}\le X_t$};
\end{tikzpicture}
\caption{鞅、下鞅、上鞅的漂移方向。三者从同一当前值 $X_t$ 出发，下一期条件期望分别上移（下鞅，对持有者有利）、持平（鞅，公平）、下移（上鞅，不利）。记忆口诀：``上鞅向下、下鞅向上'',名字来自其期望的极限性质（上鞅期望递减），与漂移方向恰好相反。}
\label{fig:martingale-2}
\end{figure}

\rmk{名字容易记反：\textbf{上}鞅的期望随时间\emph{递减}、路径系统性\emph{下}沉；\textbf{下}鞅则期望递增、向上漂。术语对的是``期望序列的单调性''（super 对应优超、势能递减），不是路径的方向，初学者务必当心。}

由塔性质，鞅的\emph{无条件}期望恒定，这是最常被用到的一条简单推论。

\prop{鞅的期望恒定}{
若 $\set{X_t}$ 是鞅，则 $\E{X_t}=\E{X_0}$ 对一切 $t$ 成立。对下鞅 $\E{X_t}$ 关于 $t$ 递增，对上鞅递减。
}
\pf{
对鞅，用塔性质把信息一层层抹掉：
\[
    \E{X_{t+1}}=\E{\E{X_{t+1}\given\cF_t}}=\E{X_t},
\]
故 $\E{X_t}$ 不随 $t$ 变，等于 $\E{X_0}$。下鞅情形把中间的等号换成 $\ge$ 得 $\E{X_{t+1}}\ge\E{X_t}$（递增），上鞅换成 $\le$（递减）。
}

\section{典型例子}
\label{sec:martingale-examples}

定义抽象，例子才让它活起来。下面三个例子分别对应``公平赌局''、``统计推断''、``资产价格''——后文的全部应用都从它们生长出来。

\ex[对称随机游走是鞅]{
设 $\set{\xi_t}$ 独立同分布，$\Prob{\xi_t=+1}=\Prob{\xi_t=-1}=\tfrac12$（公平的一步赌注），$S_0=0$，$S_t=\sum_{s=1}^{t}\xi_s$。取自然 filtration $\cF_t=\sigma(\xi_1,\dots,\xi_t)$。证明 $\set{S_t}$ 是鞅。
}
\sol{
$S_t$ 是 $\xi_1,\dots,\xi_t$ 的函数，故 $\cF_t$-可测、适应；$\abs{S_t}\le t$ 故可积。验证鞅等式：
\[
    \E{S_{t+1}\given\cF_t}=\E{S_t+\xi_{t+1}\given\cF_t}
    =S_t+\E{\xi_{t+1}\given\cF_t}=S_t+\E{\xi_{t+1}}=S_t,
\]
第二步用线性与``$S_t$ 在 $\cF_t$ 下已知可提出'',第三步用 $\xi_{t+1}$ 与 $\cF_t$ 独立（故条件期望退化为无条件），末步因 $\E{\xi_{t+1}}=0$。增量 $\xi_t$ 正是鞅差分序列。\rmk{若改成有偏赌局 $\Prob{\xi=1}=p$，则 $\E{\xi}=2p-1$：$p>\tfrac12$ 时 $S_t$ 是下鞅（向上漂），$p<\tfrac12$ 时是上鞅——赌场的优势正是把玩家财富做成了上鞅。}
}

\ex[似然比是鞅]{
设 $X_1,X_2,\dots$ 在真实密度 $f_0$ 下独立同分布，$f_1$ 是另一个备择密度（在 $f_0$ 的支撑上 $f_0>0$）。定义累积似然比
\[
    L_t=\prod_{s=1}^{t}\frac{f_1(X_s)}{f_0(X_s)},\qquad L_0=1 .
\]
取 $\cF_t=\sigma(X_1,\dots,X_t)$。证明在真实分布 $f_0$ 下 $\set{L_t}$ 是鞅。
}
\sol{
$L_t$ 适应、且 $\E{L_t}=1<\infty$（下面顺带得到）。关键一步是算单步条件期望。$L_{t+1}=L_t\cdot\dfrac{f_1(X_{t+1})}{f_0(X_{t+1})}$，其中 $L_t$ 在 $\cF_t$ 下已知，$X_{t+1}$ 与 $\cF_t$ 独立，故
\[
    \E[f_0]{L_{t+1}\given\cF_t}
    =L_t\cdot\E[f_0]{\frac{f_1(X_{t+1})}{f_0(X_{t+1})}} .
\]
而那个无条件期望（在真实密度 $f_0$ 下积分）恰好为一：
\[
    \E[f_0]{\frac{f_1(X)}{f_0(X)}}
    =\int \frac{f_1(x)}{f_0(x)}\,f_0(x)\,\d x
    =\int f_1(x)\,\d x=1 .
\]
于是 $\E[f_0]{L_{t+1}\given\cF_t}=L_t$，$\set{L_t}$ 是鞅，且 $\E{L_t}=\E{L_0}=1$。\rmk{这是序贯似然比检验（下一节``序贯检验''）与测度变换的引擎：似然比鞅在 $f_0$ 下期望恒为 $1$，正是 Wald 序贯检验的边界概率得以控制、以及金融里``从客观测度换到风险中性测度''（Radon--Nikodym 导数过程）能成立的根本原因。}
}

\ex[Doob 鞅：对固定目标的逐步逼近]{
设 $Y$ 是任一可积随机变量（``终将揭晓的真相''，如某资产的清算价值、某参数的真值），$\set{\cF_t}$ 是 filtration。定义
\[
    M_t:=\E{Y\given\cF_t} .
\]
则 $\set{M_t}$ 是鞅。
}
\sol{
$M_t$ 按构造 $\cF_t$-可测、可积。用塔性质（$\cF_t\se\cF_{t+1}$，``粗投影吃掉细投影''）：
\[
    \E{M_{t+1}\given\cF_t}=\E{\E{Y\given\cF_{t+1}}\given\cF_t}=\E{Y\given\cF_t}=M_t .
\]
故 $\set{M_t}$ 是鞅，称为 $Y$ 关于该 filtration 的 \emph{Doob 鞅}。
}

最后这个例子值得专门点出，因为它把上一章的``后验信念是鞅''一句话彻底解释清楚了。

\state{后验信念是鞅 —— 一个 Doob 鞅}{
设真相是参数 $\theta$，把待估的目标取成示性量 $Y=\indicator_{\{\theta=\theta^\ast\}}$（``状态是否为 $\theta^\ast$''）。则 $M_t=\E{Y\given\cF_t}=\Prob{\theta=\theta^\ast\given\cF_t}$ 正是\emph{$t$ 期的后验信念}。由上例，它自动是鞅：$\E{M_{t+1}\given\cF_t}=M_t$。这正是上一章``你说服不了自己往哪边走''的精确来源——理性后验信念是对一个固定真相的 Doob 鞅，其期望恒等于当前信念，未来的信念变动事前无法预测。贝叶斯劝说里的 Bayes-plausibility 约束、宏观学习里``信念不可被系统性操纵''，都是这一条 Doob 鞅性的化身。
}

\section{可选停时定理}
\label{sec:martingale-ost}

鞅的``公平''在固定的确定时刻 $t$ 上由 $\E{X_t}=\E{X_0}$ 表达。但赌徒真正想问的是：我能不能用一个\emph{聪明的停手策略}打破公平——比如``只要赢了就收手、输了就接着赌''？停手的时刻本身是随机的、取决于已发生的牌局。要回答这个问题，先得给``依赖历史的随机停手时刻''一个干净的定义。

\defn{停时}{
关于 filtration $\set{\cF_t}$，一个取值于 $\set{0,1,2,\dots}\cup\set{\infty}$ 的随机变量 $\tau$ 称为 \textbf{停时}（stopping time），若对每个 $t$，事件 $\set{\tau=t}\in\cF_t$。
}

这个条件的含义朴素而关键：\emph{是否在时刻 $t$ 停手，只能依据到 $t$ 为止的信息来决定}——你不许偷看未来再决定何时收手。``第一次财富触及 $100$ 元就停''是停时（每天都能判断今天是否首次达标）；而``在财富即将下跌的前一刻停''不是停时（需要知道明天）。这正是把``可实施的交易/退出策略''与``事后诸葛''区分开的数学线。

\thm{可选停时定理（Optional Stopping Theorem）}{
设 $\set{X_t}$ 是关于 $\set{\cF_t}$ 的鞅，$\tau$ 是停时。若下列三条充分条件之一成立：
\begin{itemize}
    \item[(i)] $\tau$ 有界（存在常数 $N$ 使 $\tau\le N$ 几乎必然）；或
    \item[(ii)] $\tau<\infty$ 几乎必然，且 $\set{X_{t\wedge\tau}}$ 一致有界（存在 $C$ 使 $\abs{X_{t\wedge\tau}}\le C$）；或
    \item[(iii)] $\E{\tau}<\infty$，且增量一致有界（存在 $C$ 使 $\abs{X_{t+1}-X_t}\le C$）；
\end{itemize}
则
\[
    \E{X_\tau}=\E{X_0} .
\]
对上鞅相应有 $\E{X_\tau}\le\E{X_0}$，对下鞅 $\E{X_\tau}\ge\E{X_0}$。
}

\rmkb[为什么需要那些技术条件：能不能``无中生有'']{
没有任何附加条件时定理会失效，反例正是著名的\emph{倍注策略}（martingale betting，术语的来源）：在对称随机游走里押注，每输一次就把赌注翻倍，直到第一次赢为止收手。令 $\tau$ 为首次赢的时刻——它几乎必然有限，而停手时净赢利恒为 $+1$，于是 $\E{X_\tau}=1\neq 0=\E{X_0}$，公平似乎被打破了！破绽在于：这个策略要求``无限的本金与无限的信用''——中途的亏损没有下界，$\set{X_{t\wedge\tau}}$ 不一致有界、$\E{\tau}$ 也未必配上有界增量，定理三条都不满足。现实中本金有限（财富有下界 $-b$），一旦撞到下界就被迫离场，公平立刻恢复。定理的技术条件，本质上是在排除``用无限资源套出确定收益''这种物理上不可能的事。
}

下面给定理在最干净的有界情形 (i) 下的证明——它只是塔性质的一次迭代，把``公平在每个固定时刻成立''推广到``公平在有界的随机时刻也成立''。

\pf[有界停时情形 (i)]{
设 $\tau\le N$。关键观察：把鞅按``停或不停''分解，停时之前每一步仍是公平的。具体地，考虑\emph{停止过程} $X_{t\wedge\tau}$（撞到 $\tau$ 后就冻结不动），它仍是鞅——因为在 $\set{\tau>t}$（尚未停）上它走一步真正的鞅增量、条件期望持平，在 $\set{\tau\le t}$（已停）上它不动、条件期望也持平。形式地，
\[
    X_{(t+1)\wedge\tau}-X_{t\wedge\tau}=\indicator_{\{\tau>t\}}\,(X_{t+1}-X_t),
\]
而 $\indicator_{\{\tau>t\}}=1-\indicator_{\{\tau\le t\}}$ 是 $\cF_t$-可测的（停时定义保证 $\set{\tau\le t}\in\cF_t$）。于是
\[
    \E{X_{(t+1)\wedge\tau}-X_{t\wedge\tau}\given\cF_t}
    =\indicator_{\{\tau>t\}}\,\E{X_{t+1}-X_t\given\cF_t}
    =\indicator_{\{\tau>t\}}\cdot 0=0 ,
\]
即 $\set{X_{t\wedge\tau}}$ 是鞅。对它取无条件期望得 $\E{X_{t\wedge\tau}}=\E{X_0}$ 对一切 $t$ 成立。取 $t=N$，因 $\tau\le N$ 故 $N\wedge\tau=\tau$，于是 $\E{X_\tau}=\E{X_0}$。其余情形 (ii)(iii) 用控制收敛/一致可积把上式中的极限 $t\to\infty$ 取进期望即可，技术细节本书从略。
}

\state{可选停时定理 = ``没有稳赚的停手策略''}{
这是鞅最具经济学分量的一条。它说：面对一个公平博弈（鞅），\emph{任何}在满足温和条件下的停手规则 $\tau$ 都不能改变期望收益——$\E{X_\tau}=\E{X_0}$，赢面与输面恰好抵消。``低买高卖''听起来稳赚，但若价格是鞅、且你不能偷看未来（$\tau$ 必须是停时），你就无法系统性地``总在高点卖出''。这把``有效市场上不存在稳赚的择时策略''翻译成了一行数学。它也是\emph{无套利}的微观机制：若某停手策略能保证 $\E{X_\tau}>\E{X_0}$，那就是一台免费印钞机，理性市场会立刻把它套利掉、使价格回到鞅。
}

图~\ref{fig:martingale-3} 把这件事画在一个有上下界的公平赌局里：无论路径如何游走、无论你设定怎样的（合规的）停手规则，停时处取值的期望都钉在初值 $X_0$ 上。

\begin{figure}[h]
\centering
\begin{tikzpicture}[scale=1.0,>=Stealth,line cap=round,line join=round]
  % 可选停时：公平赌局（鞅）从 X_0 出发，停在触及上界 +a 或下界 -b 时；
  % 无论怎样选停时规则，E[X_tau]=X_0。
  \draw[->] (-0.2,0) -- (11.2,0) node[right] {$t$};
  \draw[->] (0,-2.4) -- (0,2.9) node[above] {$X_t$};

  \def\up{2.0}
  \def\lo{-2.0}
  \def\start{0.6}

  \draw[thick,green!50!black] (0,\up) -- (10.8,\up);
  \draw[thick,red!70!black]  (0,\lo) -- (10.8,\lo);
  \node[green!50!black,anchor=south west] at (0.1,\up) {上界 $+a$};
  \node[red!70!black,anchor=north west]  at (0.1,\lo) {下界 $-b$};

  \draw[dashed,blue!55!black] (0,\start) -- (10.8,\start);
  \node[anchor=east] at (-0.1,\start) {$X_0$};

  \draw[very thick,black!75]
      (0,\start) -- (0.5,1.1) -- (1.0,0.6) -- (1.5,1.1) -- (2.0,0.6)
      -- (2.5,0.1) -- (3.0,0.6) -- (3.5,0.1) -- (4.0,-0.4) -- (4.5,0.1)
      -- (5.0,0.6) -- (5.5,1.1) -- (6.0,0.6) -- (6.5,1.1) -- (7.0,1.6)
      -- (7.5,1.1) -- (8.0,1.6) -- (8.5,1.1) -- (9.0,1.6) -- (9.5,2.0);
  \fill[green!50!black] (9.5,\up) circle (2.4pt);
  \fill[black!75] (0,\start) circle (2.2pt);

  \draw[dashed,gray!60] (9.5,0) -- (9.5,\up);
  \node[anchor=north] at (9.5,-0.05) {$\tau$};

  \node[blue!55!black,anchor=west,align=left] at (6.5,-1.45)
       {$\E{X_\tau}=X_0$};
  \draw[blue!55!black,->] (6.45,-1.35) -- (5.4,{\start-0.05});
\end{tikzpicture}
\caption{可选停时定理的图像。公平赌局（鞅）从 $X_0$ 出发，在触及上界 $+a$ 或下界 $-b$ 时停手（图中样本于 $\tau$ 处触上界被吸收）。尽管单条路径或赢或输，但停时取值的\emph{期望}恒等于初值 $X_0$（蓝色虚线）——任何合规的停手规则都无法系统性地改变它。}
\label{fig:martingale-3}
\end{figure}

\ex[赌徒破产问题：用可选停时一步算出胜率]{
对称随机游走 $S_t$ 从 $S_0=k$ 出发（$0<k<n$），在触及 $0$（破产）或 $n$（达标）时停手，停时记 $\tau$。求最终达标（先到 $n$）的概率 $p_k=\Prob{S_\tau=n}$。
}
\sol{
$\set{S_t}$ 是鞅，停时 $\tau$ 把过程限制在 $[0,n]$ 内、故 $\set{S_{t\wedge\tau}}$ 一致有界，且可证 $\tau<\infty$ 几乎必然——满足可选停时定理条件 (ii)。于是
\[
    \E{S_\tau}=\E{S_0}=k .
\]
而 $S_\tau$ 只取两个值：以概率 $p_k$ 取 $n$、以概率 $1-p_k$ 取 $0$。代入：
\[
    \E{S_\tau}=n\,p_k+0\cdot(1-p_k)=n\,p_k=k
    \quad\Longrightarrow\quad
    p_k=\frac{k}{n} .
\]
达标概率正比于起点离破产线的距离——直觉完全合理（起点越靠近目标越容易先到），而我们没有解任何差分方程，只用了``公平博弈的期望守恒''。\rmk{若把目标换成``求期望持续时间 $\E{\tau}$'',对称随机游走的另一个鞅 $S_t^2-t$（可验证它也是鞅）配可选停时定理，可一行得到 $\E{\tau}=k(n-k)$。鞅方法把这类边界问题统一成``找一个合适的鞅 + 套可选停时''。}
}

\section{鞅收敛定理}
\label{sec:martingale-convergence}

鞅不仅在每步保持公平，在温和条件下还\emph{必然收敛}到一个极限随机变量。这条定理是``学习终将稳定''、``信念终将沉淀''这类断言的数学根据。其证明（Doob 的上穿不等式）较重，本书按惯例只陈述结论并给直觉。

\thm{鞅收敛定理（Doob）}{
设 $\set{X_t}$ 是鞅（或下鞅、上鞅），且在 $L^1$ 上一致有界：
\[
    \sup_t\,\E{\abs{X_t}}<\infty .
\]
则存在一个可积的随机变量 $X_\infty$，使得 $X_t\asto X_\infty$（几乎必然收敛）。特别地，非负上鞅总满足这一条件，故\emph{非负上鞅必几乎必然收敛}。
}

直觉是这样的：鞅不能有系统性漂移（期望恒定），所以它要么逐渐安定下来、要么无界地剧烈摆动；而``$L^1$ 一致有界''恰好排除了``摆动幅度无限增大''这条出路。Doob 的证明把``不收敛''刻画为``无穷多次上穿任一给定区间 $[a,b]$'',再证明鞅的``上穿次数''期望有限（否则会与公平性矛盾），于是几乎必然只能上穿有限次——也就只能收敛。

\state{鞅收敛 = ``学习终将沉淀''}{
把这条定理用在 Doob 鞅 $M_t=\E{Y\given\cF_t}$（后验信念 / 对真相的逐步逼近）上，威力立现。后验概率 $M_t\in[0,1]$ 是有界鞅，故必几乎必然收敛到某个 $M_\infty$。这说明：理性贝叶斯学习者的信念\emph{不会永远震荡，终将沉淀}到一个极限信念。若信息流足够丰富（$\cF_\infty$ 能识别真相），则 $M_\infty$ 收敛到真相本身——这就是贝叶斯一致性的鞅证明。在宏观学习、社会学习（信息瀑布）、重复博弈的信念演化里，``信念收敛''这一步几乎总是鞅收敛定理在背后兜底。
}

\rmkb[收敛了，期望就一定守得住吗：一致可积]{
几乎必然收敛 $X_t\to X_\infty$ 并\emph{不}自动保证 $\E{X_\infty}=\E{X_0}$（期望可能在极限处``漏掉'')。回到对称随机游走加吸收：它是鞅、有界停止后收敛，但若不设上界、只在触 $0$ 时停，则 $S_\tau\equiv 0$，$\E{S_\tau}=0\neq k=\E{S_0}$——期望守恒失效。补上``一致可积''这条额外条件，才有 $\E{X_\infty}=\E{X_0}$ 且 $X_t\to X_\infty$ 在 $L^1$ 中成立。直觉上一致可积排除了``质量逃向无穷''。这与可选停时定理需要技术条件是同一回事：极限/停时与期望交换，需要``不漏质量''的保证。
}

\section{它通向何处}
\label{sec:martingale-applications}

鞅是少数几个``一处学会、四处遇见''的工具。把本章接回经济学主干，四个去处尤为重要。

\subsection{有效市场与风险中性定价}

有效市场假说最锋利的数学形式就是一句话：\textbf{（贴现后的）资产价格是鞅}。若价格 $P_t$ 关于公开信息流 $\cF_t$ 满足 $\E{P_{t+1}\given\cF_t}=P_t$（更准确地，是相对无风险利率贴现后的价格），则由可选停时定理，任何基于公开信息的择时策略都无法系统性获利——``可预测的超额收益''不存在。这正是上一章``后验信念是鞅''在价格上的体现：价格作为对未来现金流的理性信念，其变动只能是不可预测的``意外''。

更深一层是\emph{无套利与风险中性定价}。资产定价基本定理断言：市场无套利，当且仅当存在一个等价的``风险中性测度'' $\QQ$，使得每个资产的贴现价格在 $\QQ$ 下是鞅：
\[
    \frac{P_t}{(1+r)^t}=\E[\QQ]{\frac{P_{t+1}}{(1+r)^{t+1}}\given\cF_t} .
\]
于是任何衍生品的当期价格 $=$ 其未来收益在风险中性测度下的贴现期望——\emph{定价就是对一个鞅取条件期望}。从客观测度 $\PP$ 换到风险中性测度 $\QQ$，用的正是似然比鞅那样的 Radon--Nikodym 导数过程（见``似然比是鞅''例）。Black--Scholes 公式、二叉树定价、利率模型，骨架全是这一条鞅性质。

\subsection{理性预期：预测误差是鞅差分}

理性预期假设个体的主观预期等于真实的条件期望：$P^e_{t+1}=\E{P_{t+1}\given\cF_t}$。于是\emph{预测误差}
\[
    e_{t+1}:=P_{t+1}-\E{P_{t+1}\given\cF_t}
\]
满足 $\E{e_{t+1}\given\cF_t}=0$——它是一个\textbf{鞅差分序列}。这给出理性预期最可检验的推论：\emph{预测误差不可被任何 $t$ 期信息预测}，即 $e_{t+1}$ 与一切 $\cF_t$-可测变量正交（$\E{e_{t+1}\,Z_t}=0$）。计量上这就是一组可检验的正交矩条件：若回归预测误差对任何已知信息（滞后变量、其它公开数据）得到显著系数，理性预期假设即被拒绝。宏观里对``预期是否理性''的检验、对 Fed 沟通效果的评估，做的都是这件事。

\subsection{时间序列的鞅中心极限定理}

经典中心极限定理要求独立同分布，时间序列里往往不成立。鞅差分序列提供了恰到好处的替代：只要 $\set{D_t}$ 是鞅差分（$\E{D_{t+1}\given\cF_t}=0$）且满足温和的矩与稳定性条件，就有
\[
    \frac{1}{\rtn}\sum_{t=1}^{n}D_t\ \dto\ \cN(0,\sigma^2) .
\]
这就是\textbf{鞅中心极限定理}。它之所以是渐近计量的主力，是因为绝大多数估计量的``得分''（一阶条件的样本平均）在正确设定下天然是鞅差分——回忆``条件期望''一章里 M-估计量由样本一阶条件隐含定义，那个一阶条件正是鞅差分。于是 MLE、GMM、时间序列回归的渐近正态性，统统由鞅 CLT 而非 i.i.d.\ 的 CLT 来保证。鞅差分``两两不相关''（本章已证）正是让方差 $\sigma^2$ 不出现协方差交叉项、CLT 得以干净成立的关键。

\subsection{序贯检验与随机逼近}

似然比鞅（$\E{L_t}=1$）撑起\emph{序贯检验}：Wald 的序贯概率比检验（SPRT）让样本量本身成为停时——边观察边累积证据，似然比一旦穿越上/下阈值就停手判决。可选停时定理保证了这种``适应性停手''不会暗中扭曲第一类/第二类错误率，从而能在期望意义下用最少的样本完成检验。这是 A/B 测试``偷看数据''之所以危险、而正确的序贯设计之所以安全的根本原因。鞅还是\emph{随机逼近}（Robbins--Monro，随机梯度下降的祖先）收敛性证明的核心工具：把迭代误差构造成一个上鞅，再用鞅收敛定理证明它几乎必然趋零——现代机器学习优化算法的收敛性分析，许多就走这条鞅路线。

\medskip

一个由``条件期望持平''定义的简单对象，串起了金融、宏观、计量与统计推断。这再次印证全书的主线：鞅之所以能同时是``有效市场''、是``理性预期误差''、是``渐近正态的得分''、是``收敛的学习'',正因为它们共享同一个内核——\emph{给定信息，最优预测就是当下；变化只能是不可预测的意外}。认清这把钥匙，许多扇看似无关的门便一起打开了。
